% !TEX root = ../main.tex
% Abstract
%--------------------------------------
\vspace*{45pt}
\begin{flushleft}
	{\Large \textbf{\scshape{Abstract}}}
\end{flushleft}
\vspace*{10pt}

The relationship between a protein's amino acid sequence and its three-dimensional structure and biological function is a central focus of modern computational biology. This project presents the development of a comprehensive web-based platform, built with the Django framework, designed to predict both the structure and the functional characteristics of proteins directly from FASTA sequences. For functional annotation, the platform integrates the DeepGO API to predict Gene Ontology (GO) terms, offering rapid and accessible sequence-to-function mapping. To address the complex protein folding problem, the study utilizes the 2D Hydrophobic-Polar (HP) lattice model. A suite of classical heuristic optimization algorithms including Hill Climbing, Simulated Annealing, Monte Carlo, and Replica Exchange Monte Carlo was implemented and benchmarked to establish structural baselines. Building upon these classical methods, Artificial Intelligence techniques were introduced to map folding trajectories. The progression from Tabular Q-Learning to a custom PyTorch-based Deep Q-Network (DQN) demonstrates the capability of Convolutional Neural Networks (CNNs) to learn spatial dependencies within the protein grid, overcoming the dimensional limitations of tabular methods. The resulting platform unifies these diverse computational methodologies into a single interface, demonstrating the potential of Deep Reinforcement Learning for complex biophysical optimization problems.

\vspace*{20pt}

\noindent \textbf{Keywords}:
Protein Folding, Deep Reinforcement Learning, Django, Gene Ontology, HP Model.
